\documentclass[11pt,letterpaper]{article}
\usepackage{jahnel}
\runningtitle{Connected runs in a Poisson line environment}
\hypersetup{pdftitle={Connected runs in an isotropic Poisson line environment},pdfauthor={}}
\title{Connected runs in an isotropic Poisson line environment}
\author{}
\date{November 3, 2026}

\begin{document}
\maketitle
\begin{abstract}
We bound the connected runs in the continuous covering used for the greedy-path reduction of isotropic Poisson line Cox percolation. The argument includes runs assembled from arbitrarily many shorter intervals. On a fixed line, the probability that a segment of length $t$ is covered by one angular family is at most $A_\mu/(1+t)$. Independence of the two angular families gives a quadratic bound for their intersection, from which the intensity of runs of length at least $t$ is at most $C_\mu t^{-3}$ per unit arclength. After integration over the line process, the corresponding planar intensity is at most $C\mu^3t^{-3}$ for small $\mu$. A boundary correction controls runs entering a region from outside. Consequently, the longest run meeting a square of radius $n$ is eventually $O(n^{2/3+\varepsilon})$ almost surely for every $\varepsilon>0$. These estimates settle the component-length and large-scale remainder questions for this covering. The bound for the total reward collected by an arbitrary path remains open.
\end{abstract}

\section{Connected-run counts and their largest scale}

The previous calculation controlled individual intervals produced by lines that pass near a fixed line. A long connected run may also be assembled from many overlapping short intervals, so that calculation did not bound the length of a whole run. We address this issue through complete coverage of a segment, then use stationarity to count components and the Mecke formula to include lines selected from the environment.

The one-dimensional union of Poisson intervals is a classical Boolean model; see~\cite[Chapter~17]{LP17}. We give the coverage argument directly for the interval law generated by the isotropic line process. Its role here is to justify a geometric input to the existing greedy-path approach; we make no claim that the underlying one-dimensional coverage method is new.

Fix $\mu>0$. Let $\cL$ be the Poisson line process with intensity $\nu_\mu(\dd\theta\,\dd s)=(\mu/\pi)\dd\theta\dd s$, where $\ell(\theta,s)=\{x:x\cdot n_\theta=s\}$, $n_\theta=(\cos\theta,\sin\theta)$, $0\le\theta<\pi$, and $s\in\R$. We use the arclength orientation $v_\theta=(-\sin\theta,\cos\theta)$ on this line, so that its points have coordinates $s n_\theta+xv_\theta$, $x\in\R$. All lengths of subsets of a line are arclengths. Probability and expectation for $\cL$ are denoted by $\P$ and $\E$.

For a deterministic line $\ell=\ell(\theta,s)$ and another nonparallel line $\ell'$, let $J(\ell,\ell')$ be the closed interval on $\ell$ at distance at most $2$ from $\ell'$. For a configuration $\eta$ of other lines, define
\[
 U_\ell^-(\eta)=\bigcup_{\ell'\in\eta:\,\theta(\ell')<\theta}J(\ell,\ell'),
 \qquad
 U_\ell^+(\eta)=\bigcup_{\ell'\in\eta:\,\theta(\ell')>\theta}J(\ell,\ell'),
 \qquad V_\ell(\eta)=U_\ell^-(\eta)\cap U_\ell^+(\eta).
\]
Parallel pairs have zero intensity and are omitted. A run is a connected component of $V_\ell(\eta)$. For a bounded run $C$, write $|C|$ for its length and $a(C)$ for its first point in the chosen orientation. When $\ell\in\cL$, these definitions always use $\eta=\cL\setminus\{\ell\}$. Runs on different lines are counted separately, including when they are produced by some of the same lines.

For $t>0$ and a bounded Borel set $K\subset\R^2$, let $N_t(K)$ count bounded runs with $|C|\ge t$ and $a(C)\in K$. For a compact convex set $K$ with nonempty interior, let $M_t(K)$ count runs of length at least $t$ that meet $K$, including any unbounded runs. Write $|K|$ for area and $p(K)$ for perimeter. In particular, $M_t(K)$ includes components that start outside the region.

We give explicit, deliberately nonoptimal constants. For $u\ge8$ put $F_\mu(u)=\mu\e^{4\mu}(64+12\log(1+u))/u$, and let $T_\mu$ be the least integer at least $\max(8,216\mu)$ for which $F_\mu(T_\mu)\le1/6$. This integer exists because $F_\mu(u)\to0$. Define
\[
 A_\mu=
 \begin{cases}
  192\mu,&0<\mu\le1/32,\\
  T_\mu+1,&\mu>1/32,
 \end{cases}
 \qquad B_\mu=8A_\mu^2,
 \qquad C_*=8\cdot192^2.
\]
The next theorem bounds complete connected components, with their multiplicity across different lines retained.

\begin{theorem}\label{thm:runs}
For every $\mu>0$, almost surely all runs on all lines of $\cL$ are bounded. For $t\ge1$ and bounded Borel $K$,
\begin{equation}\label{eq:starts}
 \E N_t(K)\le B_\mu\mu\,\frac{|K|}{t^3}.
\end{equation}
For $t\ge1$ and compact convex $K$ with nonempty interior,
\begin{equation}\label{eq:meeting}
 \E M_t(K)\le B_\mu\mu
       \Big(\frac{|K|}{t^3}+\frac{p(K)}{\pi t^2}\Big).
\end{equation}
For $0<\mu\le1/32$, the factor $B_\mu\mu$ in both bounds equals $C_*\mu^3$.
\end{theorem}

For $n\ge1$, set $Q_n=[-n,n]^2$ and let $L_n$ be the supremum of the lengths of runs meeting $Q_n$, with supremum zero when there are none. We use $O(\cdot)$ almost surely to allow a finite random constant and a random threshold. The following consequence supplies the cutoff needed when summing an additive error over dyadic scales.

\begin{corollary}\label{cor:maximum}
For every $\mu>0$ and $\varepsilon>0$, almost surely $L_n=O(n^{2/3+\varepsilon})$ as $n\to\infty$.
\end{corollary}

\begin{proof}
Put $a=2/3+\varepsilon$ and apply \eqref{eq:meeting} with $n=2^j$ and $t=n^a$. Since $|Q_n|=4n^2$ and $p(Q_n)=8n$, Markov's inequality gives
\[
 \P(L_{2^j}\ge2^{ja})
 \le B_\mu\mu\big(4\,2^{-3\varepsilon j}
                +(8/\pi)2^{-(1/3+2\varepsilon)j}\big).
\]
The bound is summable in $j$, so the Borel--Cantelli lemma implies $L_{2^j}<2^{ja}$ eventually. For $2^j\le n<2^{j+1}$, monotonicity gives $L_n\le L_{2^{j+1}}\le2^a n^a$ once $j$ is sufficiently large. This proves the assertion for all real $n\ge1$, and hence also for integer scales.
\end{proof}

\section{Proof of the connected-run estimates}

The principal input is a coverage bound on one deterministic line. The angular restrictions defining $U_\ell^-$ and $U_\ell^+$ use disjoint parts of the Poisson line process, so their coverage probabilities multiply. Stationarity then converts complete coverage into a bound on component starts, and the Mecke formula transfers that bound to every line of the random environment.

For a deterministic line $\ell$, let $U_\ell=\bigcup_{\ell'\in\cL}J(\ell,\ell')$ use all angular directions. Let $u_\mu(t)$ be the probability that an interval of length $t$ on $\ell$ is contained in $U_\ell$; isotropy and translation invariance make this independent of the chosen line and interval. The auxiliary lemma controls all chains of intervals by controlling the union itself.

\begin{lemma}\label{lem:coverage}
For every $\mu>0$ and $t\ge0$, $u_\mu(t)\le A_\mu/(1+t)$. The same bound holds when the other lines are restricted to any deterministic measurable set of normal angles.
\end{lemma}

\begin{proof}[Proof of Theorem~\ref{thm:runs}]
Fix a deterministic line $\ell$ and identify its arclength coordinate with $\R$. The two angular restrictions of $\cL$ are independent Poisson processes. Applying Lemma~\ref{lem:coverage} to each restriction gives, for every $s\ge0$,
\begin{equation}\label{eq:double}
 v_\ell(s):=\P([0,s]\subset V_\ell(\cL))
  =\P([0,s]\subset U_\ell^-(\cL))
   \P([0,s]\subset U_\ell^+(\cL))
  \le\frac{A_\mu^2}{(1+s)^2}.
\end{equation}
This tends to zero as $s\to\infty$. An unbounded component contains a half-line starting at some rational coordinate, so a countable union and the same estimate in both directions show that all components are bounded almost surely.

Only finitely many intervals $J(\ell,\ell')$ meet a fixed compact interval of $\ell$: each such $\ell'$ meets its compact neighbourhood of radius $2$, which is met by finitely many Poisson lines. Thus, $V_\ell$ is locally a finite union of closed intervals, and its component starts are locally finite. Their expected number in a bounded interval is finite as well, because starts are among the interval endpoints, which have finite intensity. The marked process of component starts and lengths is stationary. Let $\kappa_\ell$ be its length-intensity measure, so that the expected number of starts in an interval of length $D$ with lengths in a Borel set $H\subset(0,\infty)$ is $D\kappa_\ell(H)$. Degenerate components can be omitted; they have zero probability because the interval endpoints have continuous joint locations.

Campbell's formula \cite[Section~2.2]{LP17} counts the possible starts of a component containing $[0,s]$. A component of length $b$ allows a set of starts of length $(b-s)^+$, where $x^+=\max(x,0)$. Therefore,
\begin{equation}\label{eq:component-intensity}
\begin{aligned}
 v_\ell(s)&=\int_0^\infty(b-s)^+\,\kappa_\ell(\dd b),\\
 \kappa_\ell([t,\infty))
   &\le\frac2t\,v_\ell(t/2)
    \le\frac{8A_\mu^2}{t^3}=\frac{B_\mu}{t^3},\\
 \P(0\text{ belongs to a run of length at least }t)
   &\le2v_\ell(t/2)\le\frac{B_\mu}{t^2}.
\end{aligned}
\end{equation}
For the last inequality, a component of length at least $t$ containing zero covers either $[-t/2,0]$ or $[0,t/2]$. These estimates are uniform in the normal angle of $\ell$.

Let $I$ be a compact interval of length $D$ on $\ell$. A component meeting $I$ either starts in $I$ or contains its first endpoint. The first alternative has expected count at most $B_\mu D/t^3$ for components of length at least $t$, and the second has expected count at most $B_\mu/t^2$, because at most one component contains that endpoint. Thus, the expected count of all such components meeting $I$ is at most $B_\mu(D/t^3+1/t^2)$. This is the boundary correction needed for components that start outside a region.

We now integrate over the random lines. The Mecke formula \cite[Theorem~4.1]{LP17} replaces a sum over $\ell\in\cL$ with an integral over a deterministic inserted line and leaves the other lines with their original law. For each bounded set of lines meeting a fixed disk, the probability of an unbounded run on an inserted line is zero by \eqref{eq:double}. Applying the formula on these sets and then taking a countable union over disks shows that all runs on all lines of $\cL$ are bounded almost surely.

For a bounded Borel set $K$, the component-start intensity on a fixed line gives an upper bound $B_\mu\mathcal H^1(\ell\cap K)/t^3$. For convex $K$, the intersection $\ell\cap K$ is an interval, and the preceding boundary estimate applies. Here $\mathcal H^1$ denotes arclength. Integration over displacement $s$ at a fixed normal angle is Fubini's theorem in orthogonal coordinates, while the measure of lines meeting $K$ is obtained by integrating its width. Consequently,
\begin{equation}\label{eq:mecke}
\begin{aligned}
 \int\mathcal H^1(\ell\cap K)\,\nu_\mu(\dd\ell)
   &=\mu|K|,\qquad
   \nu_\mu\{\ell:\ell\cap K\ne\varnothing\}=\mu p(K)/\pi,\\
 \E N_t(K)&\le B_\mu\mu|K|/t^3,\\
 \E M_t(K)&\le B_\mu\mu
       \big(|K|/t^3+p(K)/(\pi t^2)\big).
\end{aligned}
\end{equation}
The line-measure formula in the first row is used only for convex $K$ in its perimeter term. All integrands are nonnegative, so the Mecke and Fubini steps also apply before finiteness is established. The displayed bounds give that finiteness, and the definition of $A_\mu$ yields $B_\mu\mu=C_*\mu^3$ for $\mu\le1/32$.
\end{proof}

\section{Proof of the coverage estimate}

We first use the positive vacant fraction to obtain a small coverage probability at one sufficiently large scale. We then separate two subintervals and remove the intervals that meet both of them. Their total intensity is of order $1/t$, which gives a recursion improving the initial estimate to the required reciprocal bound.

\begin{proof}[Proof of Lemma~\ref{lem:coverage}]
Represent another line by its intersection coordinate $x\in\R$ with $\ell$ and its relative direction angle $\alpha\in(0,\pi)$. The intensity is $(\mu/\pi)\sin\alpha\dd x\dd\alpha$, and the induced interval has length $d(\alpha)=4/\sin\alpha$. Moving from the center $x$ to the left endpoint preserves Lebesgue measure for each angle. The mapping and restriction properties of Poisson processes \cite[Section~5.1]{LP17} therefore give a stationary Poisson process of intervals. Its rate of starts is $2\mu/\pi$, and the mean number covering any fixed point is $4\mu$.

For $y\ge8$, the rate of starts of intervals with length greater than $y$ is
$Q_\mu(y)=(2\mu/\pi)(1-\sqrt{1-16/y^2})\le6\mu/y^2$. The inequality follows by rationalising the square root and using $\sqrt{1-16/y^2}\ge\sqrt3/2$. Let $H_\mu(s)$ be the intensity of intervals covering both zero and $s$. Then $H_\mu(s)=\int_s^\infty Q_\mu(y)\dd y\le6\mu/s$ for $s\ge8$, and $H_\mu(s)\le4\mu$ for every $s\ge0$.

Let $p=\e^{-4\mu}$ and let $W_t$ be the Lebesgue measure of the vacant set in $[0,t]$. The Poisson void probability gives $\E W_t=pt$ and the joint vacancy probability at separation $s$ equal to $p^2\e^{H_\mu(s)}$. Complete coverage implies $W_t=0$. Thus, for $t\ge8$,
\begin{equation}\label{eq:initial}
\begin{aligned}
 u_\mu(t)&\le\frac{\operatorname{Var}(W_t)}{p^2t^2}
     =\frac{2}{t^2}\int_0^t(t-s)
                   \big(\e^{H_\mu(s)}-1\big)\dd s\\
 &\le\frac{2\e^{4\mu}}{t}\int_0^tH_\mu(s)\dd s
  \le\frac{\mu\e^{4\mu}}{t}
               \big(64+12\log(1+t)\big)=F_\mu(t).
\end{aligned}
\end{equation}
This proves that the initial scale $T_\mu$ defined before the theorem has $u_\mu(T_\mu)\le1/6$.

For the recursion, full coverage of $[0,3t]$ implies coverage of both $[0,t]$ and $[2t,3t]$. An interval meeting both subintervals covers the two points $t$ and $2t$, and the number of these shared intervals is Poisson with mean $H_\mu(t)$. After their removal, the intervals that can cover the two subintervals belong to disjoint sets of Poisson intervals. The two remaining coverage events are independent, and each has probability at most $u_\mu(t)$ by monotonicity under deletion. Hence,
\begin{equation}\label{eq:recursion}
 u_\mu(3t)\le u_\mu(t)^2+H_\mu(t)
           \le u_\mu(t)^2+\frac{6\mu}{t},\qquad t\ge8.
\end{equation}
This step concerns whole unions of intervals; it allows any number of shorter intervals in either covered subinterval.

For general $\mu$, the conditions $u_\mu(T_\mu)\le1/6$ and $T_\mu\ge216\mu$ permit induction in \eqref{eq:recursion}. For $0<\mu\le1/32$, instead start at $8$ and use $u_\mu(8)\le1-\e^{-4\mu}\le4\mu$. The two inductions give
\begin{equation}\label{eq:induction}
\begin{aligned}
 u_\mu(3^jT_\mu)&\le\tfrac16\,3^{-j}
       &&(j\ge0),\\
 u_\mu(8\cdot3^j)&\le6\mu\,3^{-j}
       &&(j\ge0,\ 0<\mu\le1/32).
\end{aligned}
\end{equation}
Indeed, the first induction follows from $1/36+6\mu/T_\mu\le1/18$. For the second, its induction step reduces to $36\mu\,3^{-j}+3/4\le2$, which holds for $\mu\le1/32$.

The function $u_\mu$ is nonincreasing. Between consecutive scales in the first bound, it is at most $T_\mu/(2t)$, so $u_\mu(t)\le(T_\mu+1)/(1+t)$ for $t\ge T_\mu$; below $T_\mu$ the latter bound follows from $u_\mu(t)\le1$. Between consecutive scales in the second bound, it is at most $144\mu/t$ for $t\ge8$, which is at most $192\mu/(1+t)$. For $t\le8$, use $u_\mu(t)\le4\mu\le192\mu/(1+t)$. This proves the asserted bound with the stated $A_\mu$. Finally, restriction to any angular set deletes intervals and can only decrease the coverage probability.
\end{proof}

\section{Consequences and the remaining path estimate}

The bounds above concern the actual isotropic line environment and the continuous covering already used for its lattice reward. They include arbitrary chains of intervals and components entering a region from outside. We now explain exactly which part of the greedy-path reduction they supply and which part still requires a separate estimate.

Let $Q(z)=z+[0,1]^2$ and let $N(z)$ count the lines meeting $Q(z)$, for $z\in\Z^2$. For a finite lattice path $\gamma$, write $R(\gamma)=\sum_{z\in\gamma}N(z)$. Let $Z(\ell)$ consist of the squares met by $\ell$ where its normal angle is neither the smallest nor the largest among the lines meeting that square. The elementary count of these interior positions gives
\[
 R(\gamma)\le2\#\gamma+
             \sum_{\ell\in\cL}\#\big(\gamma\cap Z(\ell)\big).
\]
If $z\in Z(\ell)$, then $Q(z)\cap\ell\subset V_\ell(\cL\setminus\{\ell\})$: each point of this intersection is within distance $\sqrt2<2$ of a line on each angular side. Thus, every square counted by the excess is covered by a run on the corresponding line.

Write $S_\gamma=\bigcup_{z\in\gamma}Q(z)$. A segment of length $b$ meets at most $4(1+b)$ unit squares. One way to see this uniformly in its orientation is to cover it by at most $\lfloor b\rfloor+1$ subsegments of length less than one; each has horizontal and vertical span less than one and meets at most four squares. Assign each square in $\gamma\cap Z(\ell)$ to the component containing $Q(z)\cap\ell$. Almost surely no line passes through a lattice vertex, and every such intersection is a nondegenerate segment, so this assignment is well defined. Consequently,
\begin{equation}\label{eq:reward}
 R(\gamma)\le2\#\gamma+
  4\sum_{\ell\in\cL}\ 
   \sum_{\substack{C\text{ run on }\ell\\ C\cap S_\gamma\ne\varnothing}}
                     (1+|C|).
\end{equation}
The count includes the required multiplicities over lines. It is a sufficient upper bound and may count components that do not contain any square of $Z(\ell)$ visited by the path.

For $k\ge0$, put runs of length in $[2^k,2^{k+1})$ in class $k$. Fix $0<\varepsilon<1/3$ and let $k_n$ be the largest nonnegative class represented by a run meeting $Q_{4n}$, with $k_n=-1$ if no such class occurs. Corollary~\ref{cor:maximum} gives
\[
 \sum_{k=0}^{k_n}2^k\le2L_{4n}
           =O(n^{2/3+\varepsilon})=o(n)
           \qquad\text{almost surely}.
\]
An admissible path of $n$ steps starting in $Q_n\cap\Z^2$, with increments of maximum norm at most two, has $S_\gamma\subset Q_{4n}$ for $n\ge1$. Therefore, the displayed estimate applies uniformly to the additive dyadic terms for all such paths. It supplies the large-scale remainder that had previously been asserted without proof.

The factor $t^{-3}$ in \eqref{eq:starts} makes the former density heuristic an upper bound for a precisely defined collection of components. However, a first-moment density estimate does not control the maximum of the sum in \eqref{eq:reward} over all admissible paths. Components on different lines remain dependent through the lines producing them, and the contribution of short components must also be controlled. The remaining objective is a finite almost-sure linear growth rate for that weighted component sum, or a sharper bound directly for the lattice reward. No positivity statement for the critical point intensity follows from the present note alone.

\bibliographystyle{amsalpha}
\bibliography{references}
\end{document}
